% Part: lambda-calculus
% Chapter: lambda-definability
% Section: lambda-definable-recursive

\documentclass[../../../include/open-logic-section]{subfiles}

\begin{document}

\olfileid[hi]{lam}{dfl}{ldr}
\olsection{\usetoken{S}{lambda definable} फलन पुनरावर्ती हैं}

न केवल सभी आंशिक पुनरावर्ती फलन !!{lambda definable} हैं, बल्कि विलोम भी
सत्य है। अर्थात् सभी !!{lambda definable} फलन आंशिक पुनरावर्ती हैं।

\begin{thm}
  \ollabel{thm:lambda-computable} यदि कोई आंशिक फलन $f$ !!{lambda
  definable} है, तो वह आंशिक पुनरावर्ती है।
\end{thm}

\begin{proof}
  हम प्रमाण की केवल रूपरेखा देते हैं। पहले $\lambd$-पदों का अंकगणितीकरण
  करते हैं, अर्थात् अनुक्रमों के प्रचलित अभाज्य-घात कूटन से $\lambd$-पदों
  को व्यवस्थित रूप से गोडेल संख्याएँ देते हैं। फिर आंशिक पुनरावर्ती फलन
  $\fn{normalize}(t)$ परिभाषित करते हैं, जो किसी लैम्ब्डा पद की गोडेल
  संख्या~$t$ को तर्क के रूप में लेकर, यदि उसका प्रसामान्य रूप हो तो उसकी
  गोडेल संख्या लौटाता है, और अन्यथा अपरिभाषित रहता है। फिर दो आंशिक
  पुनरावर्ती फलन $\fn{toChurch}$ और $\fn{fromChurch}$ परिभाषित करते हैं,
  जो प्राकृतिक संख्याओं को अनुरूप चर्च अंक की गोडेल संख्याओं में और उनसे
  वापस प्रतिचित्रित करते हैं।

  इन पुनरावर्ती फलनों से हम फलन~$f$ को आंशिक पुनरावर्ती फलन के रूप में
  परिभाषित कर सकते हैं। कोई $\lambd$-पद~$F$, ~$f$ को !!{lambda define}
  करता है। $f(n_1, \dots, n_k)$ संगणित करने के लिए पहले
  $\fn{toChurch}(n_i)$ से अनुरूप चर्च अंकों की गोडेल संख्याएँ प्राप्त करें
  और उन्हें $\Gn{F}$ में जोड़कर पद $F \num{n_1}\dots\num{n_k}$ की गोडेल
  संख्या पाएँ। अब इस गोडेल संख्या पर $\fn{normalize}$ उपयोग करें। यदि
  $f(n_1, \dots, n_k)$ परिभाषित है, तो $F
  \num{n_1}\dots\num{n_k}$ का प्रसामान्य रूप है (जो चर्च अंक ही होना
  चाहिए); अन्यथा उसका कोई प्रसामान्य रूप नहीं है (और इसलिए
  \[\fn{normalize}(\Gn{F\num{n_1}\dots\num{n_k}})\] अपरिभाषित है)। अंत
  में प्रसामान्य किए गए पद की गोडेल संख्या पर $\fn{fromChurch}$ उपयोग करें।
\end{proof}

\end{document}
